% nf-macros.tex — visual building blocks for "Metodi Statistici per la Neuropsicologia Forense"
% Adapted from the Time Management deck (tm-macros.tex); loads unipdwarsawstyle.sty first.

% ---------------------------------------------------------
% FONTS, ICONS, LIBRARIES
% ---------------------------------------------------------
\usepackage{unipdwarsawstyle}
\usepackage[sfdefault]{FiraSans}
\usepackage{fontawesome5}
\usetikzlibrary{arrows.meta,positioning,calc,fit,backgrounds,shapes.geometric}

% ---------------------------------------------------------
% IMAGE HEIGHT CAP — 16:9 frames are shorter than the old 4:3 ones, so images sized
% by width can run off the bottom. Every \includegraphics taller than \nfmaxheight
% is scaled down (aspect ratio kept); unsized markdown images use the same limit.
% ---------------------------------------------------------
\newlength{\nfmaxheight}
\AtBeginDocument{\setlength{\nfmaxheight}{0.74\textheight}}
\makeatletter
\newsavebox\nf@imgbox
\let\nf@includegraphics\includegraphics
\renewcommand*\includegraphics[2][]{%
  \sbox\nf@imgbox{\nf@includegraphics[#1]{#2}}%
  \ifdim\dimexpr\ht\nf@imgbox+\dp\nf@imgbox\relax>\nfmaxheight
    \resizebox{!}{\nfmaxheight}{\usebox\nf@imgbox}%
  \else\usebox\nf@imgbox\fi}
% unsized markdown images: fit within \linewidth and \nfmaxheight
\renewcommand*\pandocbounded[1]{%
  \sbox\nf@imgbox{#1}%
  \Gscale@div\@tempa{\nfmaxheight}{\dimexpr\ht\nf@imgbox+\dp\nf@imgbox\relax}%
  \Gscale@div\@tempb{\linewidth}{\wd\nf@imgbox}%
  \ifdim\@tempb\p@<\@tempa\p@\let\@tempa\@tempb\fi
  \ifdim\@tempa\p@<\p@\scalebox{\@tempa}{\usebox\nf@imgbox}\else\usebox\nf@imgbox\fi}
\makeatother

% ---------------------------------------------------------
% PALETTE — structure is UniPD crimson/gold (from the style file)
% ---------------------------------------------------------
\definecolor{nfInk}{HTML}{243447}      % dark text
\definecolor{nfPoll}{HTML}{5B3F9E}     % Wooclap
\definecolor{nfPadlet}{HTML}{B8860B}   % Padlet
\definecolor{nfGrey}{HTML}{8A94A0}

% ---------------------------------------------------------
% GLOBAL LOOK
% ---------------------------------------------------------
\hyphenpenalty=10000 \exhyphenpenalty=10000
\setbeamertemplate{headline}{}
\setbeamercolor{background canvas}{bg=white}
\setbeamercolor{normal text}{fg=nfInk}
\setbeamerfont{frametitle}{size=\large,series=\bfseries}
\setbeamertemplate{footline}{%
  \hfill{\color{nfGrey}\tiny\insertframenumber}\hspace*{4mm}\vspace*{2mm}}
\setbeamertemplate{itemize item}{\color{unipdCrimson}\raisebox{0.15ex}{\scriptsize\faCircle}}

% ---------------------------------------------------------
% ICON ROW — icon in a crimson disc + text
%   \iconrow{\faVideo}{Text}
% ---------------------------------------------------------
\newcommand{\iconrow}[2]{%
  \par\noindent
  \begin{minipage}[c]{0.9cm}\centering
    \tikz\node[circle, fill=unipdCrimson, text=white, minimum size=7mm, inner sep=0pt]{\small #1};
  \end{minipage}%
  \begin{minipage}[c]{\dimexpr\linewidth-0.9cm\relax}\small #2\end{minipage}\par\vspace{5pt}}

% ---------------------------------------------------------
% ACTIVITIES — real QR code and join code
% ---------------------------------------------------------
\newcommand{\timer}[1]{%
  \tikz[baseline=(t.base)]\node[fill=unipdCrimson, text=white, rounded corners=6pt,
    font=\scriptsize\bfseries, inner sep=3pt] (t) {\faIcon{stopwatch}\ #1};}

%   \qrbox{colour}{Figures/qr.png}{code or empty}
\newcommand{\qrbox}[3]{%
  \begin{tikzpicture}
    \node[draw=#1, line width=1.2pt, rounded corners=5pt, inner sep=5pt, fill=white] (q)
      {\includegraphics[width=3.4cm]{#2}};
    \if\relax\detokenize{#3}\relax\else
      \node[below=3pt of q, font=\small\bfseries\color{#1}] {codice: #3};
    \fi
  \end{tikzpicture}}

%   \wooclapslot{tipo}{domanda}{istruzioni}{minuti o vuoto}{Figures/qr.png}{codice}
\newcommand{\wooclapslot}[6]{%
  \begin{columns}[c]
    \begin{column}{0.56\textwidth}
      {\color{nfPoll}\small\faPoll\ \textbf{WOOCLAP}\,·\,#1}%
      \if\relax\detokenize{#4}\relax\else\hfill\timer{#4}\fi\par\vspace{8pt}
      {\LARGE\bfseries\color{nfInk}#2\par}\vspace{10pt}
      {\small\color{nfInk}#3\par}
    \end{column}
    \begin{column}{0.40\textwidth}\centering\qrbox{nfPoll}{#5}{#6}\end{column}
  \end{columns}}

%   \padletslot{tipo}{domanda}{istruzioni}{minuti o vuoto}{Figures/qr.png}
\newcommand{\padletslot}[5]{%
  \begin{columns}[c]
    \begin{column}{0.56\textwidth}
      {\color{nfPadlet}\small\faStickyNote\ \textbf{PADLET}\,·\,#1}%
      \if\relax\detokenize{#4}\relax\else\hfill\timer{#4}\fi\par\vspace{8pt}
      {\LARGE\bfseries\color{nfInk}#2\par}\vspace{10pt}
      {\small\color{nfInk}#3\par}
    \end{column}
    \begin{column}{0.40\textwidth}\centering\qrbox{nfPadlet}{#5}{}\end{column}
  \end{columns}}

% ---------------------------------------------------------
% TITLE SLIDE — use on a  ## {.plain}  frame:   \nftitleslide{2 · Cornice teorica}
% ---------------------------------------------------------
\newcommand{\nftitleslide}[1]{%
  \begin{tikzpicture}[remember picture, overlay]
    \fill[unipdCrimson] (current page.north west) rectangle ([xshift=0.58\paperwidth]current page.south west);
    \fill[unipdGold] ([xshift=0.58\paperwidth]current page.north west) rectangle ([xshift=0.585\paperwidth]current page.south west);
  \end{tikzpicture}%
  \begin{columns}[c]
    \begin{column}{0.55\textwidth}
      \color{white}
      {\small\color{unipdGold}\MakeUppercase{#1}\par}\vspace{6pt}
      {\fontsize{22}{27}\selectfont\bfseries Metodi Statistici per la Neuropsicologia Forense\par}\vspace{22pt}
      {\small\bfseries Giorgio Arcara\par}
      {\scriptsize Università di Padova\\ IRCCS San Camillo, Venezia\par}
    \end{column}
    \begin{column}{0.40\textwidth}\centering
      \includegraphics[width=0.85\linewidth]{Figures/LogoSanCamilloIRCCS_Unipd_alpha.png}\par\vspace{14pt}
      \includegraphics[height=0.55cm]{Figures/CC_license_3_0.png}
    \end{column}
  \end{columns}}

% ---------------------------------------------------------
% SECTION DIVIDER — built automatically from each "# Section" heading
% (YAML: section-titles: false, so Quarto does not add its own)
% ---------------------------------------------------------
\AtBeginSection{%
  \begin{frame}[plain]
    \begin{tikzpicture}[remember picture, overlay]
      \fill[unipdCrimson] (current page.north west) rectangle ([xshift=3mm]current page.south west);
      \fill[unipdGold] ([xshift=3mm]current page.north west) rectangle ([xshift=3.6mm]current page.south west);
      \node[anchor=north west, font=\scriptsize\color{nfGrey}, inner sep=0pt]
        at ([xshift=8mm,yshift=-5mm]current page.north west)
        {METODI STATISTICI PER LA NEUROPSICOLOGIA FORENSE};
    \end{tikzpicture}
    \vfill
    \begin{center}
      {\fontsize{26}{32}\selectfont\bfseries\color{unipdCrimson}\insertsectionhead\par}
    \end{center}
    \vfill
  \end{frame}}

% ---------------------------------------------------------
% TAKEAWAY strip at the bottom of the slide
% ---------------------------------------------------------
\newcommand{\takeaway}[1]{%
  \begin{tikzpicture}[remember picture, overlay]
    \node[anchor=south west, fill=unipdCream!70, text=nfInk, font=\scriptsize,
          inner xsep=6pt, inner ysep=3pt, rounded corners=2pt, text width=0.84\paperwidth]
      at ([xshift=5mm,yshift=4mm]current page.south west)
      {\textcolor{unipdCrimson}{\faKey}\ \textbf{In sintesi:} #1};
  \end{tikzpicture}}
